% ====================================================================
% KONFIGURASI FORMAT STRUKTUR DOKUMEN
% ====================================================================
% File ini mengatur format untuk chapter, section, subsection, dll.
% Format disesuaikan dengan Juknis UNEJ Versi 2.0 (2023)
%
% Ketentuan:
% - Chapter (BAB): 14pt, Bold, All Caps, Centered
% - Section (bagian): Bold, Normal size, Penomoran
% - Subsection & Subsubsection: Bold, Normal size
% ====================================================================

% --- KONFIGURASI CHAPTER (BAB) ---
% Format: Centered, Bold, All Caps (14pt sesuai Juknis)
% Spasi sebelum: -20pt, sesudah: 40pt
\titleformat{\chapter}[block]
  {\bfseries\fontsize{14pt}{16pt}\selectfont\centering}
  {\MakeUppercase{\chaptertitlename} \thechapter.\hspace{0.5em}}
  {0pt}
  {\MakeUppercase}

\titlespacing*{\chapter}{0pt}{-20pt}{40pt}  % Spacing: normal, before=-20pt, after=40pt

% --- KONFIGURASI SECTION (Bagian 1) ---
% Format: Bold, Normal size (12pt), Bernomor
\titleformat{\section}
  {\bfseries\normalsize}               % Style: Bold, Normal size
  {\thesection}                        % Label: "1", "1.1", dll
  {0.5em}                                % Space antara label dan title
  {}                                   % Formatting title (none)

% --- KONFIGURASI SUBSECTION (Bagian 1.1) ---
% Format: Bold, Normal size (12pt), Bernomor
\titleformat{\subsection}
  {\bfseries\normalsize}               % Style: Bold, Normal size
  {\thesubsection}                     % Label: "1.1", "1.1.1", dll
  {0.5em}                                % Space antara label dan title
  {}                                   % Formatting title (none)

% --- KONFIGURASI SUBSUBSECTION (Bagian a., b., dll) ---
% Format: Normal font (12pt, tidak bold), Format huruf kecil (a., b., dll)
% Tidak masuk ke Daftar Isi (tocdepth = 2)
\setcounter{secnumdepth}{3}
\setcounter{tocdepth}{2}
\renewcommand{\thesubsubsection}{\alph{subsubsection}.}
\titleformat{\subsubsection}[block]
  {\normalfont\normalsize}             % Style: Normal font (TIDAK BOLD), Normal size
  {\hspace*{0.5cm}\thesubsubsection}   % Label: "a.", "b.", dll (menjorok sedang 0.5cm)
  {0.5em}                              % Space antara label dan title
  {}                                   % Formatting title (none)

\titlespacing*{\subsubsection}{0pt}{10pt}{4pt}  % Spacing: indent=0pt, before=10pt, after=4pt

% --- KONFIGURASI PARAGRAPH ---
% Untuk case yang lebih spesifik (jarang digunakan dalam juknis)
\titleformat{\paragraph}[runin]
  {\bfseries}
  {\theparagraph}
  {0.5em}
  {}

% --- KONFIGURASI IDENTASI ---
\setlength{\parindent}{1cm}  % Atur jarak indentasi

% --- KONFIGURASI KETEBALAN GARIS TABEL (STANDAR JURNAL AKADEMIK) ---
% Menyeragamkan ketebalan garis pada seluruh tabel (booktabs dan standard) agar rapi, proporsional, dan elegan
\setlength{\heavyrulewidth}{0.75pt}  % Garis atas (\toprule) & garis bawah (\bottomrule)
\setlength{\lightrulewidth}{0.45pt}  % Garis pemisah header (\midrule)
\setlength{\cmidrulewidth}{0.35pt}   % Garis sub-kolom (\cmidrule)
\setlength{\arrayrulewidth}{0.45pt}  % Garis pemisah standar (\hline pada tabular/longtable)
\setlength{\aboverulesep}{0.35ex}    % Jarak vertikal di atas garis booktabs
\setlength{\belowrulesep}{0.35ex}    % Jarak vertikal di bawah garis booktabs

% --- KONFIGURASI DAFTAR ISI / TABEL / GAMBAR ---
% Dipusatkan di file ini supaya styling tetap konsisten dengan struktur dokumen
\renewcommand{\cftchapleader}{\cftdotfill{\cftdotsep}}

% --- KONFIGURASI JUDUL DAFTAR ISI / TABEL / GAMBAR ---
% Menyamakan format judul dengan \chapter: Bold, 14pt, Centered, ALL CAPS

\renewcommand{\cfttoctitlefont}{%
  \hfill\bfseries\fontsize{14pt}{16pt}\selectfont\MakeUppercase
}
\renewcommand{\cftaftertoctitle}{\hfill}

\renewcommand{\cftlottitlefont}{%
  \hfill\bfseries\fontsize{14pt}{16pt}\selectfont\MakeUppercase
}
\renewcommand{\cftafterlottitle}{\hfill}

\renewcommand{\cftloftitlefont}{%
  \hfill\bfseries\fontsize{14pt}{16pt}\selectfont\MakeUppercase
}
\renewcommand{\cftafterloftitle}{\hfill}



% Hapus after-title skip default tocloft agar konsisten
\setlength{\cftaftertoctitleskip}{40pt}
\setlength{\cftafterlottitleskip}{40pt}
\setlength{\cftafterloftitleskip}{40pt}
\setlength{\cftbeforetoctitleskip}{-20pt}
\setlength{\cftbeforelottitleskip}{-20pt}
\setlength{\cftbeforeloftitleskip}{-20pt}
